\chapter{Options: Cross-Venue and Volatility Arbitrage at Speed}\label{ch:hf:options-cross-venue-and-volatility-arbitrage-at-speed}

The day before a stock goes ex-dividend, thousands of deep in-the-money calls should be exercised and some are not; the holders who forget pay the
ones who remember, and the trade that collects it has been in the academic literature for nearly twenty years. Pool, Stoll and Whaley found
that more than half of the long positions that should have been exercised were not, over ten years to 2006, at a cost to holders of over \$491
million. In this chapter's model the play starts to pay when 1.7\% of holders forget; at 10\% it earns \$1.05 for every public contract open,
at 50\% \$10.

\section{Parity at speed: conversions and reversals}

Put--call parity (Book 1, chapter 25) ties a call, a put and the forward. When quotes violate it beyond costs, three legs lock in the difference.

\begin{definition}[Conversion arbitrage, reversal arbitrage]\label{def:hf:options-cross-venue-and-volatility-arbitrage-at-speed:conversion}
\emph{Conversion arbitrage}\index{conversion arbitrage} buys the stock, buys a put and sells a call with the same strike and expiry, locking in the
strike at expiry, when the call is rich relative to the put and the stock. \emph{Reversal arbitrage}\index{reversal arbitrage} does the opposite:
shorts the stock, sells the put and buys the call, when the call is cheap; it pays the stock's borrow and any dividends.
\end{definition}

With many venues quoting the same series, the best bid and offer of each option across venues may violate parity even when no single venue does.
In the chapter's scan (\cref{lst:hf:options-cross-venue-and-volatility-arbitrage-at-speed:scan}), nine strikes of a three-month chain are quoted ten
cents wide on each venue around fair values from Black-76, each venue's mid off by its own noise of three cents; the stock is a cent either side of
100; each leg costs a cent. A violation appears in 0.7\% of snapshots with one venue, 6.6\% with three and 26\% with six, worth one to two cents a
share. Then the quotes refresh once, keeping a share of their noise: when half persists, only 6--8\% of the violations are still there; when 90\%
persists, 43\%. The arbitrage exists in proportion to the number of venues and survives in proportion to how slowly they update.

\section{Box spreads and the rate they imply}

A box (One Quant Book 5, chapter 1) combines a bull call spread and a bear put spread with the same strikes: it pays $K_2-K_1$ at expiry whatever
happens, so its price is a zero-coupon bond. A box between 95 and 105 expiring in three months at \$9.90 lends at 4.02\% continuously compounded; if
the arbitrageur's own rate is 4\%, buying it at \$9.89 earns a cent. Van Binsbergen, Diamond and Grotteria used this: inferring risk-free rates from
option prices without a model of risk, they found the convenience yield on Treasuries equal to about 40 basis points, larger below three months and
four times larger during the financial crisis. The box's implied rate is a market rate for secured lending through the clearing house, and trading
the box against one's own funding is arbitrage only within the costs of the four legs and the capital the positions consume.

\begin{definition}[Jelly roll]\label{def:hf:options-cross-venue-and-volatility-arbitrage-at-speed:jelly}
A \emph{jelly roll}\index{jelly roll} is a long synthetic forward (long call, short put) in one expiry against a short synthetic forward in another,
at the same strike; its price is the difference of two forwards and implies the cost of carry (financing less dividends and borrow) between the two
expiries.
\end{definition}

In the chapter's chain, a \omterm{def:hf:options-cross-venue-and-volatility-arbitrage-at-speed:jelly}{jelly roll} between the three- and six-month expiries implies a carry of 4.0\%, the model's rate, since no dividend falls
between them. A dividend announced for that period, or a change in the stock's borrow, moves the \omterm{def:hf:options-cross-venue-and-volatility-arbitrage-at-speed:jelly}{jelly roll} before it moves anything else.

\section{Early exercise, dividends and the dividend play}

American calls on dividend-paying stocks should be exercised the day before the ex-date when the dividend exceeds the time value given up (One Quant
Book 5, chapter 6). For a call struck at 40 on a \$60 stock with a month left, a dividend above 13 cents makes exercise optimal. Holders who fail to
exercise leave their short counterparts unassigned: those keep a call that has lost its dividend. Assignment falls pro rata on all short open
interest, so market makers who trade large offsetting positions with each other just before the ex-date and exercise all their longs take most of
the assignments that the forgetful holders did not cause (Book 5's \texttt{firm.optmm.best\_play}). Hao, Kalay and Mayhew found this scheme
inflating reported volume on the last cum-dividend day, and facilitated by the exchanges' limits on the transaction costs passed to traders.

\begin{omfigure}
\begin{tikzpicture}
\begin{axis}[width=11cm,height=5.4cm,axis y line*=right,axis x line=none,xmode=log,log basis x=10,xmin=0.9,xmax=55,ymin=0,ymax=100,
  ylabel={\small value captured (\%)},tick label style={font=\footnotesize},table/col sep=comma]
\addplot[omDef,thick,dashed,mark=square*,mark size=1.3pt] table[x=fail,y=captured]{figdata/hft/20-options-cross-venue-and-volatility-arbitrage-at-speed/dividend.csv};
\label{plot:hf:divcap}
\end{axis}
\begin{axis}[width=11cm,height=5.4cm,axis y line*=left,xmode=log,log basis x=10,xmin=0.9,xmax=55,xtick={1,2,5,10,20,50},
  xticklabels={1,2,5,10,20,50},xlabel={\small holders failing to exercise (\%)},ylabel={\small net \$ per public contract},ymin=0,ymax=11,
  tick label style={font=\footnotesize},grid=major,grid style={gray!25},table/col sep=comma,
  legend style={font=\scriptsize,at={(0.03,0.97)},anchor=north west}]
\addplot[omThm,thick,mark=*,mark size=1.3pt] table[x=fail,y=per_public]{figdata/hft/20-options-cross-venue-and-volatility-arbitrage-at-speed/dividend.csv};
\addlegendentry{net profit (left)}
\addlegendimage{/pgfplots/refstyle=plot:hf:divcap}\addlegendentry{share captured (right)}
\end{axis}
\end{tikzpicture}

\omcaption{The dividend play at its best size: net profit of the two market makers per contract of public open interest (10,000 contracts), and
the share of the value left by failing holders that they capture, against the share of holders who fail to exercise (log scale); 30 cents a share
gained on each unassigned short call, trade fees of 0.2 cent and exercise fees of 0.1 cent a share. Data: \texttt{hf\_optarb.dividend}.}\label{fig:hf:options-cross-venue-and-volatility-arbitrage-at-speed:dividend}
\end{omfigure}

The play's profit rises with the share that fails (\cref{fig:hf:options-cross-venue-and-volatility-arbitrage-at-speed:dividend}), and so does the
share of the failures' value the players capture: 42\% at 5\% failing, 82\% at half. Below the break-even, $f<c/(2g)$ with $c$ the fees per contract
traded and $g$ the gain per unassigned call, the play does not pay at any size: 1.7\% with the chapter's fees.

\section{The same volatility on several listings}

An index's options trade on the index itself (cash-settled, European) and on its exchange-traded fund (physically settled, American, with
dividends); a stock's options trade on several exchanges, identical and fungible. The first pair should carry the same volatility after adjusting
for exercise style and dividends; the second, the same prices. A market maker quoting both hedges one with the other, and the spread between their
implied volatilities, net of the adjustments, is a relative-value trade that lives in seconds when one listing's quotes lag the other's
(chapter 8's lead--lag, applied to volatility).

\section{Strategy files}

\begin{strategyfile}{Conversion and reversal arbitrage}\label{strat:hf:options-cross-venue-and-volatility-arbitrage-at-speed:conversion}
\sfield{Who pays you, and why} Venues whose quotes lag the others', and traders who cross wide on one venue.
\sfield{Instruments and venues} A stock, its calls and puts on every options exchange.
\sfield{Signal} Parity violations of the best quotes across venues beyond three legs' costs, the borrow and dividends.
\sfield{Sizing and execution} All three legs at once, the options on the venues quoting the violation; hold to expiry or unwind when parity returns.
\sfield{Costs} Three legs' fees, the borrow for reversals, early assignment risk on American options.
\sfield{How it dies} Speed: in the model, with quotes keeping half their noise, only 6--8\% of violations survive one refresh.
\sfield{Horizon, capacity, infrastructure} Microseconds to expiry; every venue's feed.
\sfield{Backtest honestly} Quotes as received, the refresh that follows, the borrow and dividends as known then.
\sfield{Sources} Book 1, chapter 25; this chapter.
\end{strategyfile}

\begin{strategyfile}{Box-spread financing arbitrage}\label{strat:hf:options-cross-venue-and-volatility-arbitrage-at-speed:box}
\sfield{Who pays you, and why} Borrowers and lenders who use boxes to finance, and Treasury holders who pay for convenience.
\sfield{Instruments and venues} European-style index options (no early exercise), four legs.
\sfield{Signal} The box's implied rate against the firm's funding and Treasury rates.
\sfield{Sizing and execution} Lend or borrow through the box when its rate beats funding by more than the four legs' costs.
\sfield{Costs} Four legs' fees and spreads, margin, balance sheet.
\sfield{How it dies} It is a financing trade with a thin margin; the convenience yield is the market's price for Treasuries' safety, not a free
lunch.
\sfield{Horizon, capacity, infrastructure} Weeks to years; capacity from the box market's depth.
\sfield{Backtest honestly} The firm's actual funding rate and capital charges.
\sfield{Sources} van Binsbergen, Diamond and Grotteria (2022).
\end{strategyfile}

\begin{strategyfile}{Dividend play on unexercised calls}\label{strat:hf:options-cross-venue-and-volatility-arbitrage-at-speed:dividend}
\sfield{Who pays you, and why} Call holders who fail to exercise before the ex-date.
\sfield{Instruments and venues} Deep in-the-money calls on dividend-paying stocks, on the last cum-dividend day.
\sfield{Signal} Calls whose optimal exercise is certain (dividend above the time value), and the open interest.
\sfield{Sizing and execution} Trade large offsetting positions with another participant, exercise all longs; size by the expected failure share.
\sfield{Costs} Trade and exercise fees (what the exchanges' caps limit); the risk that everyone exercises.
\sfield{How it dies} Holders learning to exercise, fee changes, rules on such trades.
\sfield{Horizon, capacity, infrastructure} One day per ex-date; capacity from the failing open interest.
\sfield{Backtest honestly} Assignment by the clearing house's actual process; fees as capped at the time.
\sfield{Sources} Pool, Stoll and Whaley (2008); Hao, Kalay and Mayhew (2010).
\end{strategyfile}

\begin{strategyfile}{Index against ETF options volatility arbitrage}\label{strat:hf:options-cross-venue-and-volatility-arbitrage-at-speed:indexetf}
\sfield{Who pays you, and why} Quoters in the listing whose surface lags, and users who pay for one listing's features.
\sfield{Instruments and venues} Index options and the fund's options on the same index.
\sfield{Signal} Implied volatility difference, adjusted for exercise style, dividends and settlement.
\sfield{Sizing and execution} Quote or take the lagging listing; hedge vega in the other.
\sfield{Costs} Two listings' spreads; the adjustments' model error.
\sfield{How it dies} Early exercise and dividend surprises in the fund's options.
\sfield{Horizon, capacity, infrastructure} Seconds to days.
\sfield{Backtest honestly} Both surfaces on one clock; the adjustments as known at the time.
\sfield{Sources} Chapter 8; this chapter.
\end{strategyfile}

\section{Tutorial: the calls nobody exercised}\label{tut:hf:options-cross-venue-and-volatility-arbitrage-at-speed:tut}

\begin{tutorial}
\textbf{Goal.} Scan a chain on several venues for parity violations and measure how many survive; price a box and a \omterm{def:hf:options-cross-venue-and-volatility-arbitrage-at-speed:jelly}{jelly roll}; size the dividend
play. \textbf{End state:} \cref{fig:hf:options-cross-venue-and-volatility-arbitrage-at-speed:dividend} and the numbers in the text.
\begin{tutsteps}
\item \textbf{The scan} of conversions and reversals across venues with \texttt{firm.parity}.
\omcode{code/firm/optarb/firm_optarb.py}{79}{109}{Best quotes across venues, three legs' fees, and a second look after one refresh.}\label{lst:hf:options-cross-venue-and-volatility-arbitrage-at-speed:scan}
\item \textbf{Boxes and \omterm{def:hf:options-cross-venue-and-volatility-arbitrage-at-speed:jelly}{jelly rolls}}.
\omcode{code/firm/optarb/firm_optarb.py}{38}{52}{The box's implied rate and edges; the carry implied by two synthetic forwards.}\label{lst:hf:options-cross-venue-and-volatility-arbitrage-at-speed:box}
\item \textbf{The dividend play} for each share of failing holders (\texttt{hf\_optarb.dividend}); the exercise threshold with
\texttt{firm.american}.
\end{tutsteps}
\textbf{What to change next.} Add the borrow fee to reversals and see which violations remain; replace the noise by quotes that lag the stock by a
venue-specific delay; estimate the failure share from open interest before and after past ex-dates.
\end{tutorial}

\section{Build: the options arbitrage module}\label{bld:hf:options-cross-venue-and-volatility-arbitrage-at-speed:optarb}

\begin{build}
\bfield{Purpose} Scan listed options for parity, box and carry violations across venues, and size the dividend play.
\bfield{Interface} \texttt{box\_rate}, \texttt{box\_edge}, \texttt{jelly\_roll}, \texttt{Listings(n\_venues, strikes, seed)} with \texttt{refresh}
and \texttt{best}, \texttt{scan(listings, spot\_bid, spot\_ask, fee, lag\_rho, rounds)}, \texttt{dividend\_curve(fails)},
\texttt{exercise\_threshold}. Built on \texttt{firm.parity} (Book 1), \texttt{firm.american} and \texttt{firm.optmm} (Book 5).
\bfield{Rules} Dollars per share; fees per leg; the dividend play's gain and fees per share, times 100 a contract.
\bfield{Acceptance tests} \texttt{code/firm/optarb/tests/}: box rate and edges by hand; the \omterm{def:hf:options-cross-venue-and-volatility-arbitrage-at-speed:jelly}{jelly roll} recovers the carry; no violation without noise;
more persistent noise lets more violations survive; the play pays nothing below its break-even and more as more holders fail.
\bfield{Stretch} Borrow and hard-to-borrow stocks; venue delays; SPX against SPY volatility.
\end{build}

\begin{omsources}
\item R.~Pool, H.~R.~Stoll, R.~E.~Whaley, Failure to exercise call options: an anomaly and a trading game, \emph{Journal of Financial Markets}
11(1), 2008, 1--35.
\item J.~Hao, A.~Kalay, S.~Mayhew, Ex-dividend arbitrage in option markets, \emph{Review of Financial Studies} 23(1), 2010, 271--303.
\item J.~H.~van Binsbergen, W.~F.~Diamond, M.~Grotteria, Risk-free interest rates, \emph{Journal of Financial Economics} 143(1), 2022, 1--29.
\end{omsources}

\section{Exercises}

\begin{exercise}[$\star$]\label{exo:hf:options-cross-venue-and-volatility-arbitrage-at-speed:1}
A 95--105 box expiring in three months trades at \$9.90. What rate does it imply?
\end{exercise}

\begin{exercise}[$\star$]\label{exo:hf:options-cross-venue-and-volatility-arbitrage-at-speed:2}
Call bid 4.60, put ask 4.10, stock ask 100.01, strike 100, three months, rate 4\%, dividends worth 0.50 before expiry. What is the conversion's edge?
\end{exercise}

\begin{exercise}[$\star$]\label{exo:hf:options-cross-venue-and-volatility-arbitrage-at-speed:3}
With trade fees of 0.2 cent and exercise fees of 0.1 cent a share and a gain of 30 cents per unassigned call, what failure share makes the dividend
play pay?
\end{exercise}

\begin{exercise}[$\star\star$]\label{exo:hf:options-cross-venue-and-volatility-arbitrage-at-speed:4}
Why do more venues produce more parity violations, and why do they not produce more profit?
\end{exercise}

\begin{exercise}[$\star\star$]\label{exo:hf:options-cross-venue-and-volatility-arbitrage-at-speed:5}
Why is a box's implied rate above the Treasury rate of the same maturity?
\end{exercise}

\begin{exercise}[$\star\star$]\label{exo:hf:options-cross-venue-and-volatility-arbitrage-at-speed:6}
What does the dividend play do to the options' reported volume, and why does it matter to others?
\end{exercise}

\begin{exercise}[$\star\star\star$]\label{exo:hf:options-cross-venue-and-volatility-arbitrage-at-speed:7}
\emph{Coding.} With \texttt{firm.american.dividend\_threshold}, find the dividend above which a call struck at 40 on a \$60 stock with a month left
should be exercised, at 25\% volatility and 4\%.
\end{exercise}

\begin{exercise}[$\star\star\star$]\label{exo:hf:options-cross-venue-and-volatility-arbitrage-at-speed:8}
\emph{Find the flaw.} ``Our scanner found 257 parity violations per thousand snapshots across six venues, a cent each: a steady income.''
\end{exercise}

\section{Problem: The Calls Nobody Exercised}

\begin{problem}[{Weekend problem --- the calls nobody exercised}]\label{pb:hf:options-cross-venue-and-volatility-arbitrage-at-speed:1}
An options arbitrage desk runs parity scans, box trades and the dividend play.

\medskip\noindent\textbf{Part I --- Parity.}
\begin{enumerate}
\item Define conversion and \omterm{def:hf:options-cross-venue-and-volatility-arbitrage-at-speed:conversion}{reversal arbitrage}.
\item Describe the scan and its violation rates by number of venues.
\item How many survive a refresh, and what decides it?
\item Why must reversals include the borrow?
\end{enumerate}

\medskip\noindent\textbf{Part II --- Rates and carry.}
\begin{enumerate}[resume]
\item What rate does the example box imply, and what does buying it at \$9.89 earn?
\item What did van Binsbergen, Diamond and Grotteria find?
\item Define a \omterm{def:hf:options-cross-venue-and-volatility-arbitrage-at-speed:jelly}{jelly roll} and its value in the chapter's chain.
\item What moves a \omterm{def:hf:options-cross-venue-and-volatility-arbitrage-at-speed:jelly}{jelly roll} first?
\end{enumerate}

\medskip\noindent\textbf{Part III --- The dividend play.}
\begin{enumerate}[resume]
\item When is exercising a call before the ex-date optimal?
\item What did Pool, Stoll and Whaley find?
\item How does the play work, and what did Hao, Kalay and Mayhew add?
\item Give the break-even failure share.
\end{enumerate}

\medskip\noindent\textbf{Part IV --- The verdict.}
\begin{enumerate}[resume]
\item State the \emph{named result}: the dividend play's expected profit per contract as a function of the share of holders who fail to exercise.
\item Who loses in the dividend play, and could they avoid it?
\item What would lower fee caps do to the play?
\item Why are cross-listing volatility trades lead--lag trades?
\item Which strategy file carries the most capital?
\item How would you estimate the failure share before an ex-date?
\item What does a brokerage owe clients holding calls before an ex-date?
\item In one sentence: who pays for the calls nobody exercised?
\end{enumerate}
\end{problem}

\section{Interview questions}

\begin{interviewq}[$\star$ \iqroles{trader}]\label{iq:hf:options-cross-venue-and-volatility-arbitrage-at-speed:1}
Write put--call parity for a dividend-paying stock and say what you would trade if the call were a dime rich.
\end{interviewq}

\begin{interviewq}[$\star\star$ \iqroles{researcher}]\label{iq:hf:options-cross-venue-and-volatility-arbitrage-at-speed:2}
How would you measure the convenience yield of Treasuries from option prices?
\end{interviewq}

\begin{interviewq}[$\star\star$ \iqroles{trader}]\label{iq:hf:options-cross-venue-and-volatility-arbitrage-at-speed:3}
Tomorrow a stock goes ex-dividend 50 cents; you are short 1,000 deep in-the-money calls. What do you expect overnight, and what do you do?
\end{interviewq}

\begin{interviewq}[$\star\star$ \iqroles{developer}]\label{iq:hf:options-cross-venue-and-volatility-arbitrage-at-speed:4}
Design a parity scanner over 5,000 series on 16 venues that reports only executable violations.
\end{interviewq}

\begin{interviewq}[$\star\star$ \iqroles{risk}]\label{iq:hf:options-cross-venue-and-volatility-arbitrage-at-speed:5}
A reversal desk is short 2 million shares of a stock that becomes hard to borrow. What happens?
\end{interviewq}

\begin{interviewq}[$\star\star\star$ \iqroles{researcher}]\label{iq:hf:options-cross-venue-and-volatility-arbitrage-at-speed:6}
Derive the dividend play's optimal size $q^\ast=\bigl(\sqrt{aP/c}-P\bigr)/2$ from its expected profit $aq/(P+2q)-cq$.
\end{interviewq}
